% --- Archon physics formalization source begin ---
% archon:physics
% archon:covers PhyXMiniProblems/problem_phyx_mini_0877.lean
% archon:source-report reports/phyx_mini/problem_phyx_mini_0877.source.json

\chapter{Physics problem phyx\_mini\_0877}
\label{ch:PhyXMiniProblems_problem_phyx_mini_0877}

\paragraph{Problem source.}
\#\# Question

What is the equivalent capacitance of the three capacitors?

\#\# Answer choices

- A: A: 2.5\textbackslash{} \textbackslash{}mu\textbackslash{}mathrm\{F\}

- B: B: 3.5\textbackslash{} \textbackslash{}mu\textbackslash{}mathrm\{F\}

- C: C: 5.5\textbackslash{} \textbackslash{}mu\textbackslash{}mathrm\{F\}

- D: D: 7.5\textbackslash{} \textbackslash{}mu\textbackslash{}mathrm\{F\}

\#\# Figure caption (auxiliary; use the image as primary evidence)

The image is a schematic diagram of an electrical circuit. It includes the following components:

1. **Battery**: On the left, represented by two parallel lines of differing lengths (indicating positive and negative terminals).

2. **Capacitors**: 
   - There are three capacitors in the circuit.
   - The top capacitor has a capacitance of \textbackslash{}(10 \textbackslash{}, \textbackslash{}mu\textbackslash{}text\{F\}\textbackslash{}).
   - Below it, there are two capacitors in parallel:
     - The first parallel capacitor has a capacitance of \textbackslash{}(20 \textbackslash{}, \textbackslash{}mu\textbackslash{}text\{F\}\textbackslash{}).
     - The second parallel capacitor has a capacitance of \textbackslash{}(10 \textbackslash{}, \textbackslash{}mu\textbackslash{}text\{F\}\textbackslash{}).

3. **Connections**: 
   - The capacitors are connected in both series and parallel arrangements.
   - The \textbackslash{}(10 \textbackslash{}, \textbackslash{}mu\textbackslash{}text\{F\}\textbackslash{}) capacitor at the top is in series with the parallel arrangement of the \textbackslash{}(20 \textbackslash{}, \textbackslash{}mu\textbackslash{}text\{F\}\textbackslash{}) and \textbackslash{}(10 \textbackslash{}, \textbackslash{}mu\textbackslash{}text\{F\}\textbackslash{}) capacitors.

4. **Text**:
   - Each capacitor is labeled with its capacitance value next to it.

This is a typical configuration used to analyze the combination of series and parallel capacitors in circuits.

\#\# Dataset metadata

Electric Circuits; Physical Model Grounding Reasoning, Multi-Formula Reasoning

\paragraph{Recorded answer/context.}
D: D: 7.5\textbackslash{} \textbackslash{}mu\textbackslash{}mathrm\{F\}

\paragraph{Figure/image path.}
/root/proposal\_for\_physic/science-mango/phyx\_mini\_run/phyx\_data/test\_image/877.png

\paragraph{Formalization target.}
create a compiling Lean file with sorry bodies at `PhyXMiniProblems/problem\_phyx\_mini\_0877.lean`. The Lean declarations must preserve the physical quantities, dimensions or dimensional roles, figure labels, governing-law hypotheses, and final relation expressed by this problem.
Use Mathlib/Physlib names found through LeanExplore where available. If a physics API is missing, introduce faithful local abstractions rather than scalar placeholder aliases.

\begin{theorem}[Physics formalization target]
\label{thm:physics:phyx_mini_0877:target}
\lean{PhyXMiniProblems.ProblemPhyXMini0877.problem_phyx_mini_0877}
\uses{def:physics:phyx-mini-0877:phyxminiproblems-problemphyxmini0877-capacitanceinmicrofarads, def:physics:phyx-mini-0877:phyxminiproblems-problemphyxmini0877-threecapacitorcircuitsetup, def:physics:phyx-mini-0877:phyxminiproblems-problemphyxmini0877-matchesidealthreecapacitorscenario, def:physics:phyx-mini-0877:phyxminiproblems-problemphyxmini0877-matchesseriesparalleltopology, def:physics:phyx-mini-0877:phyxminiproblems-problemphyxmini0877-matchessuppliedthreecapacitorfigure, def:physics:phyx-mini-0877:phyxminiproblems-problemphyxmini0877-haspositivecircuitcapacitances, def:physics:phyx-mini-0877:phyxminiproblems-problemphyxmini0877-satisfiesidealcapacitorcombinationlaws, def:physics:phyx-mini-0877:phyxminiproblems-problemphyxmini0877-recordeddatasetanswer, def:physics:phyx-mini-0877:phyxminiproblems-problemphyxmini0877-answermatchesequivalentcapacitance}
The assigned autoformalize agent should translate this physics problem into Lean declarations in the covered file, with theorem and lemma proof bodies written as `by sorry`.
\end{theorem}
\begin{proof}
This is an autoformalization task, not a proof task. Produce statements that can later be proved by the physics prover without weakening the physical model.
\end{proof}

% --- Archon declaration topology begin ---
\section{Declaration topology}
\begin{definition}[potential Difference Dimension]
  \label{def:physics:phyx-mini-0877:phyxminiproblems-problemphyxmini0877-potentialdifferencedimension}
  \lean{PhyXMiniProblems.ProblemPhyXMini0877.potentialDifferenceDimension}
  The dimension M L² T⁻² C⁻¹ of electric potential difference
\end{definition}
\begin{proof}
  This is the declared model definition of potential Difference Dimension.
\end{proof}
\begin{definition}[capacitance Dimension]
  \label{def:physics:phyx-mini-0877:phyxminiproblems-problemphyxmini0877-capacitancedimension}
  \lean{PhyXMiniProblems.ProblemPhyXMini0877.capacitanceDimension}
  \uses{def:physics:phyx-mini-0877:phyxminiproblems-problemphyxmini0877-potentialdifferencedimension}
  The dimension C / V of electrical capacitance
\end{definition}
\begin{proof}
  This is the declared model definition of capacitance Dimension.
\end{proof}
\begin{definition}[Capacitance Quantity]
  \label{def:physics:phyx-mini-0877:phyxminiproblems-problemphyxmini0877-capacitancequantity}
  \lean{PhyXMiniProblems.ProblemPhyXMini0877.CapacitanceQuantity}
  \uses{def:physics:phyx-mini-0877:phyxminiproblems-problemphyxmini0877-capacitancedimension}
  A nonnegative, unit-independent physical electrical capacitance
\end{definition}
\begin{proof}
  This is the declared model definition of Capacitance Quantity.
\end{proof}
\begin{definition}[capacitance Readout]
  \label{def:physics:phyx-mini-0877:phyxminiproblems-problemphyxmini0877-capacitancereadout}
  \lean{PhyXMiniProblems.ProblemPhyXMini0877.capacitanceReadout}
  \uses{def:physics:phyx-mini-0877:phyxminiproblems-problemphyxmini0877-capacitancequantity}
  Read a capacitance in the coherent unit induced by a choice of base units
\end{definition}
\begin{proof}
  This is the declared model definition of capacitance Readout.
\end{proof}
\begin{definition}[capacitance In Farads]
  \label{def:physics:phyx-mini-0877:phyxminiproblems-problemphyxmini0877-capacitanceinfarads}
  \lean{PhyXMiniProblems.ProblemPhyXMini0877.capacitanceInFarads}
  \uses{def:physics:phyx-mini-0877:phyxminiproblems-problemphyxmini0877-capacitancequantity, def:physics:phyx-mini-0877:phyxminiproblems-problemphyxmini0877-capacitancereadout}
  Read a capacitance in coherent-SI farads
\end{definition}
\begin{proof}
  This is the declared model definition of capacitance In Farads.
\end{proof}
\begin{definition}[capacitance In Microfarads]
  \label{def:physics:phyx-mini-0877:phyxminiproblems-problemphyxmini0877-capacitanceinmicrofarads}
  \lean{PhyXMiniProblems.ProblemPhyXMini0877.capacitanceInMicrofarads}
  \uses{def:physics:phyx-mini-0877:phyxminiproblems-problemphyxmini0877-capacitancequantity, def:physics:phyx-mini-0877:phyxminiproblems-problemphyxmini0877-capacitanceinfarads}
  Read a capacitance in microfarads
\end{definition}
\begin{proof}
  This is the declared model definition of capacitance In Microfarads.
\end{proof}
\begin{definition}[Capacitor Label]
  \label{def:physics:phyx-mini-0877:phyxminiproblems-problemphyxmini0877-capacitorlabel}
  \lean{PhyXMiniProblems.ProblemPhyXMini0877.CapacitorLabel}
  The three capacitor symbols, named by their positions and printed values
\end{definition}
\begin{proof}
  This is the declared model definition of Capacitor Label.
\end{proof}
\begin{definition}[Circuit Node]
  \label{def:physics:phyx-mini-0877:phyxminiproblems-problemphyxmini0877-circuitnode}
  \lean{PhyXMiniProblems.ProblemPhyXMini0877.CircuitNode}
  The three electrically distinct nodes in the depicted ideal circuit
\end{definition}
\begin{proof}
  This is the declared model definition of Circuit Node.
\end{proof}
\begin{definition}[Battery Terminal]
  \label{def:physics:phyx-mini-0877:phyxminiproblems-problemphyxmini0877-batteryterminal}
  \lean{PhyXMiniProblems.ProblemPhyXMini0877.BatteryTerminal}
  The two marked terminals of the battery in the raster
\end{definition}
\begin{proof}
  This is the declared model definition of Battery Terminal.
\end{proof}
\begin{definition}[Capacitor Model]
  \label{def:physics:phyx-mini-0877:phyxminiproblems-problemphyxmini0877-capacitormodel}
  \lean{PhyXMiniProblems.ProblemPhyXMini0877.CapacitorModel}
  The component idealization used by the combination laws
\end{definition}
\begin{proof}
  This is the declared model definition of Capacitor Model.
\end{proof}
\begin{definition}[Voltage Source Model]
  \label{def:physics:phyx-mini-0877:phyxminiproblems-problemphyxmini0877-voltagesourcemodel}
  \lean{PhyXMiniProblems.ProblemPhyXMini0877.VoltageSourceModel}
  The source component shown on the left of the circuit
\end{definition}
\begin{proof}
  This is the declared model definition of Voltage Source Model.
\end{proof}
\begin{definition}[Three Capacitor Circuit Figure]
  \label{def:physics:phyx-mini-0877:phyxminiproblems-problemphyxmini0877-threecapacitorcircuitfigure}
  \lean{PhyXMiniProblems.ProblemPhyXMini0877.ThreeCapacitorCircuitFigure}
  \uses{def:physics:phyx-mini-0877:phyxminiproblems-problemphyxmini0877-capacitorlabel}
  Literal presentation information from 877.png. Printed capacitances are scalar microfarad readouts and are kept distinct from physical capacitance quantities
\end{definition}
\begin{proof}
  This is the declared model definition of Three Capacitor Circuit Figure.
\end{proof}
\begin{definition}[Three Capacitor Circuit Setup]
  \label{def:physics:phyx-mini-0877:phyxminiproblems-problemphyxmini0877-threecapacitorcircuitsetup}
  \lean{PhyXMiniProblems.ProblemPhyXMini0877.ThreeCapacitorCircuitSetup}
  \uses{def:physics:phyx-mini-0877:phyxminiproblems-problemphyxmini0877-capacitancequantity, def:physics:phyx-mini-0877:phyxminiproblems-problemphyxmini0877-capacitorlabel, def:physics:phyx-mini-0877:phyxminiproblems-problemphyxmini0877-circuitnode, def:physics:phyx-mini-0877:phyxminiproblems-problemphyxmini0877-batteryterminal, def:physics:phyx-mini-0877:phyxminiproblems-problemphyxmini0877-capacitormodel, def:physics:phyx-mini-0877:phyxminiproblems-problemphyxmini0877-voltagesourcemodel, def:physics:phyx-mini-0877:phyxminiproblems-problemphyxmini0877-threecapacitorcircuitfigure}
  Independent physical data for the circuit. Neither equivalent capacitance is defined from the component values or from the recorded answer
\end{definition}
\begin{proof}
  This is the declared model definition of Three Capacitor Circuit Setup.
\end{proof}
\begin{definition}[Matches Ideal Three Capacitor Scenario]
  \label{def:physics:phyx-mini-0877:phyxminiproblems-problemphyxmini0877-matchesidealthreecapacitorscenario}
  \lean{PhyXMiniProblems.ProblemPhyXMini0877.MatchesIdealThreeCapacitorScenario}
  \uses{def:physics:phyx-mini-0877:phyxminiproblems-problemphyxmini0877-threecapacitorcircuitsetup}
  The ideal battery-and-capacitor interpretation stated by the problem
\end{definition}
\begin{proof}
  This is the declared model definition of Matches Ideal Three Capacitor Scenario.
\end{proof}
\begin{definition}[Matches Series Parallel Topology]
  \label{def:physics:phyx-mini-0877:phyxminiproblems-problemphyxmini0877-matchesseriesparalleltopology}
  \lean{PhyXMiniProblems.ProblemPhyXMini0877.MatchesSeriesParallelTopology}
  \uses{def:physics:phyx-mini-0877:phyxminiproblems-problemphyxmini0877-threecapacitorcircuitsetup}
  The node identifications read from the primary raster. The two lower capacitors share both terminals and hence form a parallel branch; the top capacitor lies in series between that branch and the battery's positive node
\end{definition}
\begin{proof}
  This is the declared model definition of Matches Series Parallel Topology.
\end{proof}
\begin{definition}[Matches Supplied Three Capacitor Figure]
  \label{def:physics:phyx-mini-0877:phyxminiproblems-problemphyxmini0877-matchessuppliedthreecapacitorfigure}
  \lean{PhyXMiniProblems.ProblemPhyXMini0877.MatchesSuppliedThreeCapacitorFigure}
  \uses{def:physics:phyx-mini-0877:phyxminiproblems-problemphyxmini0877-capacitanceinmicrofarads, def:physics:phyx-mini-0877:phyxminiproblems-problemphyxmini0877-threecapacitorcircuitsetup}
  Primary-raster evidence and calibration of its printed microfarad labels to the independent physical component capacitances. This structure contains no equivalent-capacitance value
\end{definition}
\begin{proof}
  This is the declared model definition of Matches Supplied Three Capacitor Figure.
\end{proof}
\begin{definition}[Has Positive Circuit Capacitances]
  \label{def:physics:phyx-mini-0877:phyxminiproblems-problemphyxmini0877-haspositivecircuitcapacitances}
  \lean{PhyXMiniProblems.ProblemPhyXMini0877.HasPositiveCircuitCapacitances}
  \uses{def:physics:phyx-mini-0877:phyxminiproblems-problemphyxmini0877-capacitanceinfarads, def:physics:phyx-mini-0877:phyxminiproblems-problemphyxmini0877-threecapacitorcircuitsetup}
  Nondegeneracy of every physical capacitance used in reciprocal laws
\end{definition}
\begin{proof}
  This is the declared model definition of Has Positive Circuit Capacitances.
\end{proof}
\begin{definition}[Satisfies Ideal Capacitor Combination Laws]
  \label{def:physics:phyx-mini-0877:phyxminiproblems-problemphyxmini0877-satisfiesidealcapacitorcombinationlaws}
  \lean{PhyXMiniProblems.ProblemPhyXMini0877.SatisfiesIdealCapacitorCombinationLaws}
  \uses{def:physics:phyx-mini-0877:phyxminiproblems-problemphyxmini0877-capacitancereadout, def:physics:phyx-mini-0877:phyxminiproblems-problemphyxmini0877-threecapacitorcircuitsetup}
  The ideal-capacitor combination laws, stated for every coherent unit system: parallel capacitances add, while reciprocal capacitances add in series. These are governing relations and contain none of the requested numerical result or answer-choice data
\end{definition}
\begin{proof}
  This is the declared model definition of Satisfies Ideal Capacitor Combination Laws.
\end{proof}
\begin{definition}[Answer Choice]
  \label{def:physics:phyx-mini-0877:phyxminiproblems-problemphyxmini0877-answerchoice}
  \lean{PhyXMiniProblems.ProblemPhyXMini0877.AnswerChoice}
  The four answer labels supplied with the problem
\end{definition}
\begin{proof}
  This is the declared model definition of Answer Choice.
\end{proof}
\begin{definition}[displayed Capacitance In Microfarads]
  \label{def:physics:phyx-mini-0877:phyxminiproblems-problemphyxmini0877-answerchoice-displayedcapacitanceinmicrofarads}
  \lean{PhyXMiniProblems.ProblemPhyXMini0877.AnswerChoice.displayedCapacitanceInMicrofarads}
  \uses{def:physics:phyx-mini-0877:phyxminiproblems-problemphyxmini0877-answerchoice}
  The equivalent capacitance printed for an answer choice, in microfarads
\end{definition}
\begin{proof}
  This is the declared model definition of displayed Capacitance In Microfarads.
\end{proof}
\begin{definition}[recorded Dataset Answer]
  \label{def:physics:phyx-mini-0877:phyxminiproblems-problemphyxmini0877-recordeddatasetanswer}
  \lean{PhyXMiniProblems.ProblemPhyXMini0877.recordedDatasetAnswer}
  \uses{def:physics:phyx-mini-0877:phyxminiproblems-problemphyxmini0877-answerchoice}
  The answer label recorded in the source dataset
\end{definition}
\begin{proof}
  This is the declared model definition of recorded Dataset Answer.
\end{proof}
\begin{definition}[Answer Matches Equivalent Capacitance]
  \label{def:physics:phyx-mini-0877:phyxminiproblems-problemphyxmini0877-answermatchesequivalentcapacitance}
  \lean{PhyXMiniProblems.ProblemPhyXMini0877.AnswerMatchesEquivalentCapacitance}
  \uses{def:physics:phyx-mini-0877:phyxminiproblems-problemphyxmini0877-capacitanceinmicrofarads, def:physics:phyx-mini-0877:phyxminiproblems-problemphyxmini0877-threecapacitorcircuitsetup, def:physics:phyx-mini-0877:phyxminiproblems-problemphyxmini0877-answerchoice}
  A displayed choice agrees with the circuit's physical equivalent capacitance
\end{definition}
\begin{proof}
  This is the declared model definition of Answer Matches Equivalent Capacitance.
\end{proof}
% --- Archon declaration topology end ---
% --- Archon physics formalization source end ---
